\section{Sequential compactness}
\begin{definition}\label{def:equicontinuity}
A family $\mathcal F$ in $C([0,1],\mathbb R)$ is equicontinuous if for
every $\varepsilon>0$ there is $\delta>0$ such that $|x-y|<\delta$ implies
$|f(x)-f(y)|<\varepsilon$ for all $x,y\in[0,1]$ and $f\in\mathcal F$.
\end{definition}
\begin{theorem}[Arzel\`a--Ascoli criterion]\label{thm:criterion}
A family $\mathcal F\subset C([0,1],\mathbb R)$ is relatively compact
in the supremum norm if and only if it is uniformly bounded and
equicontinuous. Uniform boundedness means that one constant $M$ satisfies
$|f(x)|\le M$ for every $f\in\mathcal F$ and every $x\in[0,1]$.
\end{theorem}
In particular, every sequence from such a family has a subsequence
converging uniformly to a continuous real-valued function.
This last sentence describes the sequential form of the same result.
